\begin{prob}
    Consider the code below:

    \begin{minted}[autogobble]{python}
        def foo(n):
            i = 1
            while i**2 < n:
                i += 1
            return i
    \end{minted}

    \begin{subprobset}
        \begin{subprob}
            What does \python{foo(n)} compute, roughly speaking?

            \begin{soln}
                It computes, approximately, $n^{1/2}$. Of course, \python{foo}
                always returns an integer, so the result of the function is usually
                not exactly correct.

                More precisely, \python{foo} returns the smallest integer greater
                than or equal to the square root $n^{1/2}$. For example, $(10)^{1/2}$ is between 3 and 4;
                \python{foo(10)} returns 4.
            \end{soln}
        \end{subprob}

        \begin{subprob}
            What is the asymptotic time complexity of \python{foo}?

            You do not need to show your work for this problem, but providing
            an explanation might help earn partial credit in case your answer
            is incorrect.

            \begin{soln}
                $\Theta(n^{1/2})$
            \end{soln}
        \end{subprob}
    \end{subprobset}

\end{prob}
